\bookbackchapter[译序]{PREFACE}{译序}

或许可以说，不提到 \mbox{Magica Quartet}，就难以写成一部完整的新世纪日本动画史；
不看《魔法少女小圆》，也很难理解这部动画何以能越出一季的篇幅，成为此后十余年间仍被反复观看、争论和阐释的作品。
它当然继承了魔法少女动画的形貌：
变身、契约、愿望、战斗，以及少女之间的友爱；
然而这些人们熟悉的形式，一经纳入它的叙事，便逐渐显露出另一副面目。
愿望并不天然通向幸福，善意也不能保证善果，希望与绝望更不是彼此隔绝的两端。
全剧正是在这种转化之中，把一个看似熟悉的类型推到了它自身的边界。

《魔圆》的成就首先来自各创作部门严密的配合。
虚渊玄的剧本规定契约、愿望与代价之间的因果；
新房昭之的演出有意延宕关键事实的披露；
苍树梅的角色原案保留了魔法少女类型亲切明朗的外观，魔女结界的拼贴美术却不断破坏这种外观；
梶浦由记的音乐又给悲剧的进程以庄严的声调。
这些因素在十二集的篇幅中互相制约，共同完成了作品的叙事，不能割裂开来评说。
由此造成的影响，也不能用“黑暗”二字概括，《魔圆》使魔法少女的契约、使命和救济都成为必须重新解释的问题。
此后论及这一类型的变化，便很难绕而行。
它在动画史上的地位，便在这里。

Jen A. Blue 的《绝望即壤》所要处理的，正是这些并不容易用一两句感想说清的问题。
这不是一部剧情梗概，也不是把作品中的象征逐项列出的观剧手册。
作者以炼金术意义上的“腐化”为枢纽，把灵魂宝石、悲叹之种、魔女与圆环之理置于同一转化过程之中；
又借助后现代理论、佛教思想、女性主义、《浮士德》以及瓦茨拉夫·哈维尔的文字，考察希望如何生成绝望，绝望又如何成为希望赖以萌发的土壤。
全书的长处，不在于替《魔圆》给出一个排他的“标准答案”，而在于把散落于画面、台词、叙事结构和跨文本引用中的线索组织成一条可以追问、也可以反驳的解释道路。

译者最初观阅完《魔圆》，便陷入长考：
愿望和诅咒何以具有同一根源？
小圆的救济为何以自我消除为代价？
《叛逆》中的反抗又为何同时含有背叛与爱？
在网络探寻的过程中，偶然发现了原书以及山川贤一先生著，“因果的使者”前辈所译的《所谓成熟之牢笼》，遂译者一发不可收拾地开始了圆学之旅。
虽然作者的解释未必处处可从，却有一个不得不提及的长处：
她把观剧所得的零散印象落实为可以据作品检验的论点。
然而考察原书在当时约莫 2020 年的中文资料情况，翻译之意便油然而生。
只因原书术语繁复，征引又广，虽早有初译之念，却迟迟没有动手。
即使伏案半载后，又因译者的惰性，便将其束之高阁了。
在《瓦尔普吉斯的回天》的定档消息公布以后，译者便重新燃起奋斗之意，开始整理旧稿，这项搁置已久的工作才得以真正启动。
可不料电影一而再，再而三的延期，这也使得译者可以堂而皇之地反复拖延。
不过好在，翻译、校对工作在历时近两年后，终于，在新作上映前完成了拙译，也算是了却心愿。

翻译这部书的困难，并不在于某一种单独的语言。
动画中的日语、作者写作所用的英语、炼金术与哲学的术语、佛教概念以及欧美批评传统的表达，常在同一段落中彼此牵连。
若一味追求字面对等，译文很容易僵硬，若只求中文顺畅，又可能把原文刻意保留的歧义和推论磨平。
这里采取的办法，是先求论旨可辨，再求文气可读：
当涉及作品专名时，尽量采用中文观众已经熟悉的译法，其中大多都参考自百度百科、维基百科、Bangumi 等知名网站数据库；
当涉及学术概念时，译者一般先多处查证其在文章准确所指，以通行译名为先，若确需必要，则保留原文。
原作者的补充说明亦出现在参考文献中，译者一并将其译出，并归入脚注。
若语境或译法确有说明必要时，则另加译注。
至于章题所引哈维尔文字，译者力求兼顾出处、上下文与本章论旨，则把原文同拙译一起，在题记中贴出，并对深奥、不明所以然之处加之以译者按，供读者参考。

以上所述只是处理译文的办法，并不能据此保证译文完美准确。
原书所涉的炼金术、哲学、宗教与批评理论各有来历，而译者并未细看文后的参考文献，任何一处理解失当，都可能改变原句的分寸。
动画评论则另有一层困难：
《魔圆》的意义由画面、表演、声音和叙事时间共同构成，文字只能分析这些因素，绝无可能代替观众对作品的直接经验。
因此，本译本无意以中文之便利遮蔽英文原书之瑕，也不把作者的解释奉为定论。
它的用处不过是使无法直接阅读原书的读者得以通览全篇论证，同时提供一份可以逐条复核的中文文本罢了。
正如英国哲学家路德维希·维特根斯坦所言：“我的命题应当是以如下方式来起阐明作用的：
任何理解我的人，当他用这些命题为梯级而超越了它们时，就会终于认识到它们是无意义的。
（可以说，在登上高处之后他必须把梯子扔掉。）
他必须超越这些命题，然后他就会正确看待世界。”
读者若由此返回动画、漫画或英文原著，重新感受、判断书中的论述是否成立，译介的目的也就达到了。

在本译本以前，\mbox{homuhomu}与流风翥云已先后从事本书的中文翻译工作。
译者在翻译和校订时，均有取二家译文参校；
遇有疑难句和专名异译，尤其得益于这种比勘以及二家的译注。
前辈不仅提供了可资参考的处理，也使译者重新核查若干自以为无疑之处。
谨在此向两位译者一并致谢。
凡有所得，不敢掠美；
本译本所有错译、疏漏和不当之处，皆盼读者不吝指教，并由译者负责。

\begin{flushright}
  \BookTranslator\\
  二〇二六年八月
  
\end{flushright}
